%% Example CV - R. Ford Jackson
%% Comprehensive overview of all competencies, experience, and achievements.
%% Use as a master reference when tailoring role-specific CVs.
%%
%% Compile with: lualatex main_example.tex
%% (pdflatex often fails on modern MiKTeX with fontawesome5 font-expansion errors.)

\documentclass[11pt,a4paper,sans]{moderncv}
\moderncvstyle{banking}
\moderncvcolor{blue}

% Force the name and section headings to render in moderncv blue (color1).
% Default banking leaves them black: moderncvstylebanking.sty's \colorlet
% copies (not aliases) the pre-scheme accent colour, so the name colours are
% frozen before \moderncvcolor runs. Re-let them after. \namefont is the hook
% every name-style macro routes through, so this also works on moderncv 2.3.1
% (Debian/Ubuntu apt), which has no \firstnamestyle/\lastnamestyle at all.
\renewcommand*{\namefont}{\fontsize{34}{36}\bfseries\upshape}
\colorlet{firstnamecolor}{color1}
\colorlet{lastnamecolor}{color1}
\colorlet{namecolor}{color1}
\renewcommand*{\sectionstyle}[1]{{\sectionfont\color{color1}#1}}

\usepackage[utf8]{inputenc}
% pdflatex fallback only (the documented engine is lualatex, which skips this
% branch). Without T1 font encoding pdflatex builds accented letters with
% \accent, and the PDF text layer stores them decomposed - `e` + U+0300 rather
% than U+00E8 - so an ATS keyword match on "Genève" fails while the page looks
% right. moderncv 2.5 loads T1 itself under pdflatex; 2.3.1 (Debian/Ubuntu apt)
% does not. \ifpdftex comes from iftex, which every moderncv version loads.
\ifpdftex\usepackage[T1]{fontenc}\fi
% moderncv loads hyperref itself in an \AtEndPreamble hook, so \hypersetup
% must go in an \AtEndPreamble of our own: on moderncv < 2.4 a top-level
% \usepackage{hyperref} clashes with the class's own
% \RequirePackage[unicode]{hyperref}. From 2.4.0 the class passes its options
% through \PassOptionsToPackage instead, which is what removes that clash.
\AtEndPreamble{\hypersetup{
    colorlinks=true,
    linkcolor=blue,
    filecolor=magenta,
    urlcolor=blue,
    pdftitle={R. Ford Jackson - CV},
    % Keep pdfpagemode=UseNone: this block runs after moderncv's own
    % \AtEndPreamble (moderncv.cls sets pdfpagemode there), so a FullScreen
    % value here would win and open every CV in fullscreen presentation mode.
    pdfpagemode=UseNone,
}}
\usepackage[scale=0.80]{geometry}
\usepackage{import}

% personal data
\name{R. Ford}{Jackson}
\address{Dallas, TX}{}{}
\phone[mobile]{404-797-7913}
\email{Ford\_Jackson@icloud.com}
\extrainfo{\href{https://www.linkedin.com/in/ford-jackson973a40156/}{LinkedIn}}

\begin{document}

\makecvtitle

% ============================================================
%     PROFILE STATEMENT
% ============================================================

\vspace{6pt}
\small{Software engineer with 3+ years of hands-on experience building full-stack web applications, from React/Knockout.js frontends to Java and serverless backends. Combines strong fundamentals in data structures, algorithms, and Agile development with practical experience delivering AWS Lambda-based features, custom GraphQL resolvers, message queue integrations (RabbitMQ, AWS SQS), and internal data-curation tools. Thrives on cross-functional collaboration and delivering creative, elegant solutions to complex technical challenges.}

% ============================================================
%     CORE COMPETENCIES
% ============================================================

\section{Core Competencies}
\vspace{1pt}
\begin{itemize}

\item \textbf{Frontend}: React, AngularJS, Knockout.js, JavaScript, TypeScript

\item \textbf{Backend}: Java, Node.js, Python, C/C++, C\#, SQL, GraphQL API design, message queues (RabbitMQ, AWS SQS)

\item \textbf{Cloud \& Serverless}: AWS, AWS Lambda, Docker, production deployment and monitoring

\item \textbf{Data Structures \& Algorithms}: AVL trees, hash tables, inverted indices, algorithm design

\item \textbf{Tools \& Practices}: Postman, Jira, Agile/Scrum, Waterfall, MVC architecture, multi-platform development (desktop and mobile)

\end{itemize}

% ============================================================
%     PROFESSIONAL EXPERIENCE
% ============================================================

\section{Professional Experience}
\vspace{3pt}
\begin{itemize}

% --- Most Recent Role ---
\item{\cventry{2025--Present}{Technical Consultant}{Pariveda Solutions}{Dallas, TX}{}{\vspace{1pt}
\begin{itemize}
    \item Contribute as a software engineer on an agile-based team working with a React/Knockout.js frontend and Java backend, with additional Node.js, SQL, TypeScript, and RabbitMQ message queue work
    \item Work as part of a large team with clear communication, engaging with stakeholders regularly to ensure deliverables meet requirements
\end{itemize}}}

\vspace{3pt}

% --- Previous Role ---
\item{\cventry{10/2022--1/2023 \& 7/2023--12/2024}{Software Engineer}{The Contract Network}{}{}{\vspace{1pt}
\begin{itemize}
    \item Built and deployed AWS serverless Lambda functions pertaining to AI functionality within the platform, using Docker for containerization and AWS SQS for asynchronous message processing
    \item Delivered custom GraphQL resolvers to the backend stack to better handle requests from users
    \item Developed internal tools to support analyst curation and monitoring of large data sets
    \item Created backend unit tests, delivered repository-wide logging, participated in code reviews, and supported production deployment/monitoring and debugging
\end{itemize}}}

\vspace{3pt}

% --- Earlier Role ---
\item{\cventry{2022--2022}{Software Engineering Intern}{AT\&T}{}{}{\vspace{1pt}
\begin{itemize}
    \item Implemented features for a web app that allowed field technicians to improve schedule management, using Node.js and SQL
    \item Added features such as universal components and updated the API to current standards
\end{itemize}}}

\end{itemize}

% ============================================================
%     EDUCATION
% ============================================================

\section{Education}
\vspace{1pt}
\begin{itemize}

\item{\cventry{2019--2023}{BA in Computer Science}{Southern Methodist University, Bobby B. Lyle School of Engineering}{Dallas, TX}{}{\vspace{1pt}
GPA: 3.79. Key topics: Algorithms, Data Structures, Agile/Waterfall Development, Scrum, MVC Architecture.
}}

\end{itemize}

% ============================================================
%     LANGUAGES
% ============================================================

\section{Languages}
\vspace{1pt}
\begin{itemize}
\item English (native).
\end{itemize}

% ============================================================
%     PUBLICATIONS (optional - omit section if none)
%     HONORS AND AWARDS (optional - omit section if none)
%     None on file as of this setup.
% ============================================================

%     HONORS AND AWARDS (optional)
% ============================================================

\section{Honors and Awards}
\vspace{1pt}
\begin{itemize}
\item{\textbf{[Award Name]} -- [Event/Organization] ([Year]).}
\end{itemize}

=======
% ============================================================
%     REFERENCES
% ============================================================

\section{References}
\vspace{1pt}
\begin{itemize}
\item{Available upon request.}
\end{itemize}

\end{document}
